\documentclass[12pt,oneside]{book}

% ==========================================================
% METADATOS
% ==========================================================
\newcommand{\universidad}{Universidad de Buenos Aires}
\newcommand{\facultad}{Facultad de Derecho}
\newcommand{\materia}{Derecho de la Integración}
\newcommand{\titulo}{Insolvencia transfronteriza en los espacios de integración: Tratados de Montevideo, Código Bustamante y Unión Europea}
\newcommand{\autor}{Isaac Edgar Camacho Ocampo}
\newcommand{\anio}{2026}

% ==========================================================
% PAQUETES
% ==========================================================
\usepackage[utf8]{inputenc}
\usepackage[T1]{fontenc}
\usepackage[spanish,es-nodecimaldot]{babel}

\usepackage[
    a4paper,
    top=2.5cm,
    bottom=2.5cm,
    left=3cm,
    right=2.5cm,
    headheight=15pt
]{geometry}

\usepackage{amsmath}
\usepackage{amssymb}
\usepackage{mathtools}

\usepackage{graphicx}
\usepackage{float}
\usepackage{caption}
\usepackage{subcaption}

\usepackage{booktabs}
\usepackage{array}
\usepackage{longtable}
\usepackage{tabularx}

\usepackage{enumitem}

\usepackage{xcolor}
\usepackage[most]{tcolorbox}

\usepackage{fancyhdr}
\usepackage{ragged2e}

\usepackage[
    colorlinks=true,
    linkcolor=blue!50!black,
    citecolor=green!40!black,
    urlcolor=blue!60!black,
    pdftitle={\titulo},
    pdfauthor={\autor},
    pdfsubject={Apuntes universitarios},
    bookmarks=true
]{hyperref}

% ==========================================================
% CONFIGURACIÓN
% ==========================================================
\graphicspath{
    {./img/}
    {./graficos/}
    {./figuras/}
}

\setcounter{tocdepth}{2}

\setlength{\parindent}{0pt}
\setlength{\parskip}{0.5em}

\setlist[itemize]{
    topsep=0.3em,
    itemsep=0.2em
}

\setlist[enumerate]{
    topsep=0.3em,
    itemsep=0.2em
}

\clubpenalty=10000
\widowpenalty=10000

\renewcommand{\arraystretch}{1.25}

% ==========================================================
% COMANDOS
% ==========================================================
\newcommand{\abs}[1]{\left|#1\right|}
\newcommand{\norm}[1]{\left\lVert#1\right\rVert}
\newcommand{\vect}[1]{\mathbf{#1}}
\newcommand{\deriv}[2]{\frac{d#1}{d#2}}
\newcommand{\pderiv}[2]{\frac{\partial#1}{\partial#2}}

% ==========================================================
% ENTORNOS / CAJAS
% ==========================================================
\newtcolorbox{definition}[1]{
    enhanced,
    breakable,
    title={Definición: #1},
    fonttitle=\bfseries,
    colback=blue!3,
    colframe=blue!50!black,
    boxrule=0.8pt,
    arc=2mm
}

\newtcolorbox{important}{
    enhanced,
    breakable,
    title={Importante},
    fonttitle=\bfseries,
    colback=yellow!8,
    colframe=orange!70!black,
    boxrule=0.8pt,
    arc=2mm
}

\newtcolorbox{example}[1]{
    enhanced,
    breakable,
    title={Ejemplo: #1},
    fonttitle=\bfseries,
    colback=green!4,
    colframe=green!40!black,
    boxrule=0.8pt,
    arc=2mm
}

\newtcolorbox{formula}[1]{
    enhanced,
    breakable,
    title={Fórmula / Regla: #1},
    fonttitle=\bfseries,
    colback=purple!4,
    colframe=purple!50!black,
    boxrule=0.8pt,
    arc=2mm
}

\newtcolorbox{exercise}[1]{
    enhanced,
    breakable,
    title={Ejercicio: #1},
    fonttitle=\bfseries,
    colback=gray!5,
    colframe=gray!60!black,
    boxrule=0.8pt,
    arc=2mm
}

\newtcolorbox{note}{
    enhanced,
    breakable,
    title={Nota},
    fonttitle=\bfseries,
    colback=white,
    colframe=gray!50,
    boxrule=0.6pt,
    arc=2mm
}

% ==========================================================
% CONFIGURACIÓN FINAL
% ==========================================================
\pagestyle{fancy}
\fancyhf{}
\fancyhead[L]{\nouppercase{\leftmark}}
\fancyhead[R]{\materia}
\fancyfoot[C]{\thepage}
\renewcommand{\headrulewidth}{0.4pt}
\renewcommand{\footrulewidth}{0pt}

\fancypagestyle{plain}{
    \fancyhf{}
    \fancyfoot[C]{\thepage}
    \renewcommand{\headrulewidth}{0pt}
}

% ==========================================================
% DOCUMENTO
% ==========================================================
\begin{document}

% ----------------------------------------------------------
% PORTADA
% ----------------------------------------------------------
\begin{titlepage}

\thispagestyle{empty}

\begin{center}

\vspace*{1cm}

{\large \textsc{\universidad}}

\vspace{0.2cm}

{\large \textsc{\facultad}}

\vspace{3cm}

{\Huge \textbf{\materia}}

\vspace{1cm}

{\LARGE \titulo}

\vspace{0.8cm}

\textit{Comprender cada concepto es el primer paso para dominarlo.}

\vspace{1.2cm}

\rule{0.70\textwidth}{0.5pt}

\vspace{0.5cm}

{\large Apuntes de estudio}

\vspace{0.5cm}

\rule{0.70\textwidth}{0.5pt}

\vfill

{\large \autor}

\vspace{0.3cm}

{\large \anio}

\end{center}

\vfill

\begin{flushleft}
\textit{Este es un modesto aporte a los estudiantes de ``\materia'' no pretende sustituir el material oficial de la materia.}
\end{flushleft}

\end{titlepage}

% ----------------------------------------------------------
% ÍNDICE
% ----------------------------------------------------------
\tableofcontents
\clearpage

% ==========================================================
% CONTENIDO
% ==========================================================

% ----------------------------------------------------------
\chapter[Presentación del tema]{Presentación del tema: la insolvencia que cruza fronteras}
% ----------------------------------------------------------

\section{Planteo general}

El tema aborda qué ocurre cuando una empresa o un comerciante deja de pagar y su patrimonio, sus acreedores o su actividad cruzan fronteras. El desarrollo parte de los fundamentos (para qué existe un régimen de insolvencia y quién controla el procedimiento), analiza el caso Swift como ejemplo de grupo empresario, compara los dos grandes criterios para decidir qué juez y qué ley se aplican (el extraterritorial y el territorialista) y, finalmente, estudia cómo lo resuelven hoy los Tratados de Montevideo y el Código Bustamante en el Mercosur, y cómo lo hace la Unión Europea mediante su reglamento.

Los contenidos se organizan en los siguientes bloques:

\begin{itemize}
    \item \textbf{Fundamentos:} castigo, prevención y protección del mercado; concurso, quiebra y salvataje.
    \item \textbf{Caso Swift:} multinacional de composición celular y descorrimiento del velo societario.
    \item \textbf{Criterios internacionales:} sistema extraterritorial (universal) y sistema territorialista.
    \item \textbf{Mercosur:} Tratados de Montevideo (1889 y 1940) y Código Bustamante; juicio único o juicios autónomos.
    \item \textbf{Derecho comparado:} Ley Modelo y Reglamento europeo; centro de intereses principales y establecimiento.
\end{itemize}

\section{Un caso hipotético para seguir el hilo del tema}

\begin{example}{«Alfajores del Plata S. A.»}
Una fábrica de alfajores artesanales tiene su casa central en Buenos Aires, un local propio con personal y caja propia en Montevideo, y clientes que compran por internet desde Brasil y Paraguay. Un día deja de pagar sus deudas. Cada bloque del tema responde preguntas sobre ese mismo caso:
\begin{itemize}
    \item \textbf{Fundamentos:} ¿qué busca el Estado: castigar al dueño, darle una oportunidad o proteger empleos, impuestos y precios?
    \item \textbf{Swift:} ¿y si la fábrica fuera solo una pieza de un grupo de empresas manejado por una única cabeza que se beneficia por otro lado?
    \item \textbf{Criterios:} ¿se hace un solo juicio para todos sus bienes y deudas (criterio extraterritorial) o uno por país (criterio territorial)?
    \item \textbf{Mercosur:} ¿qué dicen los tratados que unen a la Argentina, Uruguay, Paraguay y Brasil? ¿Importa que el local de Montevideo sea independiente?
    \item \textbf{Europa:} ¿cómo lo resolvería la Unión Europea con su «centro de intereses principales» y sus «establecimientos»?
\end{itemize}
\end{example}

% ----------------------------------------------------------
\chapter[Fundamentos de la insolvencia]{Fundamentos de la insolvencia y rol del Estado}
% ----------------------------------------------------------

\section{Regímenes distintos entre Estados}

Quien estudia el tema en la carrera ve, en general, un solo régimen: el de la República Argentina. Sin embargo, es necesario saber que existen distintas posibilidades y distintos regímenes. Esta diversidad adquiere relevancia al abordar el plano internacional, porque repercute en el sistema que cada Estado haya adoptado.

\section{De la quiebra como castigo a la etapa concursal}

En general, la insolvencia es el determinante de que se caiga en \textbf{estado de cesación de pagos}. Los sistemas jurídicos pueden adoptar criterios distintos para definir cuándo se configura (por ejemplo, un balance negativo, o los tres últimos balances negativos).
% TODO: contenido incompleto o ambiguo en las notas originales (las notas formulan los criterios como preguntas, sin indicar cuál se adopta)

\subsection{La quiebra vista como castigo}

Una ley puede adoptar un sistema que procure sacar rápidamente del mercado a la empresa o al comerciante en crisis, para evitar que los efectos nocivos se propaguen. Se concibe entonces la quiebra como un castigo: se liquida, se intenta pagar lo que se pueda a los acreedores y se aparta al sujeto que está provocando un perjuicio en el mercado, de modo que no pueda operar.

\subsection{La etapa concursal}

Con la evolución de las leyes apareció un período en que se advirtió que la cesación de pagos no siempre obedece a la culpa o a la mala fe del comerciante o del empresario: pueden surgir situaciones fuera de su control (se mencionó como ejemplo la pandemia). El ordenamiento jurídico considera entonces que no es justo castigarlo por algo sobre lo que no tenía control, y se incorpora la \textbf{etapa concursal}.

Cuando se habla de concurso se alude normalmente al \textbf{concurso preventivo}, porque esa etapa tiene por fin prevenir que se caiga en quiebra, es decir, que se liquide. Se implementan distintas herramientas:

\begin{itemize}
    \item el concordato;
    \item el acuerdo preventivo extrajudicial;
    \item los acuerdos de quita y espera.
\end{itemize}

Todos ellos procuran dar una oportunidad a la empresa para que su economía retome el buen rumbo de sus negocios y empiece a cubrir sus deudas.

\section{El mercado como centro de la mirada: efectos de la quiebra y salvataje}

Algunos sistemas jurídicos no miran tanto la relación acreedores-deudor, sino \textbf{el mercado}, y se preguntan qué consecuencias trae que una empresa entre en estado de cesación de pagos. Las hipótesis analizadas son las siguientes:

\begin{itemize}
    \item Si el deudor no paga, tampoco podrá pagar él a sus propios acreedores.
    \item El Estado, como ente recaudador, deja de cobrar impuestos (Ingresos Brutos, Ganancias, IVA, entre otros).
    \item Al mismo tiempo, el Estado tiene más gastos: debe pagar el subsidio por desempleo, prácticamente por un año, a cada uno de los empleados de la empresa quebrada.
    \item Desaparece del mercado un bien o un servicio; si hay menos oferta y la demanda se mantiene, \textbf{suben los precios}.
\end{itemize}

Para evitar estos efectos nocivos se abre una instancia más, destinada a que la empresa no caiga en la insolvencia y en la quiebra: el salvataje.

\begin{definition}{Salvataje (cramdown)}
Mecanismo en el cual se \textbf{desinteresa al propio deudor}: se realiza un cálculo aproximado de cuánto valdría la empresa, se lo desinteresa, y la empresa puede ser adquirida por cualquiera. De este modo, otro empresario con mayor capacidad, credibilidad y flujos para invertir puede sacar a flote la empresa en crisis.
\end{definition}

Estas cuestiones ya no interesan solamente a los acreedores y al deudor: el cierre de la empresa provoca en el mercado efectos nocivos que se procura evitar. Se verifica entonces si el tercero que adquiere la empresa lo hace con algún propósito irregular o porque realmente ve posibilidades de sacarla adelante y producir. Si la empresa sigue en funcionamiento, no hay desempleo, no baja el consumo de la población y continúa el pago de Ingresos Brutos, Ganancias e IVA.

Los sistemas jurídicos pueden adoptar, en su ley interna, cualquiera de estas posturas.

\begin{important}
\textbf{Tres finalidades posibles de un régimen de insolvencia}
\begin{itemize}
    \item \textbf{Liquidar rápido} (la quiebra vista como castigo): se aparta del mercado al sujeto que provoca el perjuicio.
    \item \textbf{Prevenir} (etapa concursal): concurso preventivo, concordato, acuerdo preventivo extrajudicial, acuerdos de quita y espera.
    \item \textbf{Proteger el mercado} (salvataje o cramdown): se desinteresa al deudor y un tercero con capacidad puede adquirir la empresa y sacarla a flote.
\end{itemize}
Cada ley interna de los Estados puede colocarse en cualquiera de estos puntos.
\end{important}

\section{Finalidad del procedimiento y carácter privatista o publicista}

Además de la finalidad del procedimiento (liquidar rápido, prevenir mediante un concurso preventivo o habilitar un salvataje), influye \textbf{con qué características se regula} el procedimiento. Independientemente de su finalidad, a un procedimiento judicial se le puede dar un tinte o característica \textbf{privatista} o \textbf{publicista}. Al leer la ley se infiere el tinte que el Estado pudo haberle dado.

\subsection{Tinte privatista}

\begin{itemize}
    \item Durante la etapa concursal, los acuerdos de quita y espera los negocia directamente el deudor con sus acreedores, y la operatoria que el deudor continúa realizando la controla una junta de acreedores.
    \item En el salvataje, el nuevo dueño de la empresa negocia directamente los acuerdos con los acreedores, y el primero que los obtiene se queda con la empresa. El juez y el fiscal no tienen injerencia en el resultado.
    \item El acuerdo lo aprueba la junta de acreedores, con la homologación; pero la homologación solo verifica que se cumplieron las conformidades que exige la ley (por ejemplo, los porcentajes). \textbf{El juez no puede analizar la sustancia del acuerdo}: no puede entrar a valorar si el plan de pago es lógico o no.
\end{itemize}

\subsection{Tinte publicista}

\begin{itemize}
    \item El control de si los acuerdos están bien o mal lo realiza el juez.
    \item La autorización al deudor para endeudarse durante el proceso concursal no la otorga la junta de acreedores, sino que se da con un informe del síndico.
    \item Si el síndico informa que el plan es inviable, el juez no aprueba el concordato con esos acuerdos.
\end{itemize}

\begin{longtable}{>{\raggedright\arraybackslash}p{0.22\textwidth} >{\raggedright\arraybackslash}p{0.34\textwidth} >{\raggedright\arraybackslash}p{0.34\textwidth}}
\caption{Tinte privatista y tinte publicista del procedimiento}\\
\toprule
\textbf{Aspecto} & \textbf{Tinte privatista} & \textbf{Tinte publicista} \\
\midrule
\endfirsthead
\toprule
\textbf{Aspecto} & \textbf{Tinte privatista} & \textbf{Tinte publicista} \\
\midrule
\endhead
\bottomrule
\endlastfoot
\textbf{Acuerdos de quita y espera} & Los negocia directamente el deudor con sus acreedores. & Se someten al control del Estado (juez e informe del síndico). \\
\textbf{Control de la operatoria durante el concurso} & Lo ejerce una junta de acreedores. & Lo ejerce el juez, con apoyo del síndico. \\
\textbf{Autorización para endeudarse durante el proceso concursal} & La da la junta de acreedores. & Se da con un informe del síndico. \\
\textbf{Salvataje (cramdown)} & El nuevo dueño negocia directamente con los acreedores; el primero que obtiene los acuerdos se queda con la empresa. El juez y el fiscal no tienen injerencia en el resultado. & El control de si los acuerdos están bien o mal lo hace el juez. \\
\textbf{Rol del juez ante el acuerdo} & Homologa: solo verifica que se cumplieron las conformidades de la ley; no analiza la sustancia del acuerdo. & Si el síndico informa que el plan es inviable, el juez no aprueba el concordato. \\
\end{longtable}

\subsection*{Cuadro Cornell del capítulo}
\addcontentsline{toc}{section}{Cuadro Cornell del capítulo}

\begin{longtable}{>{\raggedright\arraybackslash}p{0.27\textwidth} >{\raggedright\arraybackslash}p{0.65\textwidth}}
\toprule
\textbf{Preguntas / palabras clave} & \textbf{Notas / respuestas} \\
\midrule
\endfirsthead
\toprule
\textbf{Preguntas / palabras clave} & \textbf{Notas / respuestas} \\
\midrule
\endhead
\bottomrule
\endlastfoot

\textbf{¿Qué papel cumple la cesación de pagos en la insolvencia?} &
La insolvencia es el determinante de que se caiga en estado de cesación de pagos. Cada sistema jurídico puede adoptar criterios distintos para definirla (por ejemplo, un balance negativo o los tres últimos balances negativos), y de ello depende cómo se abre el procedimiento. \\

\textbf{¿Cuáles son las tres finalidades posibles de un régimen de insolvencia?} &
Liquidar rápido (la quiebra como castigo), prevenir (concurso preventivo, concordato, acuerdo preventivo extrajudicial, quita y espera) y proteger el mercado (salvataje o cramdown). Cada ley interna puede ubicarse en cualquiera de ellas. \\

\textbf{¿Por qué se incorporó la etapa concursal?} &
Porque se advirtió que la cesación de pagos no siempre se debe a culpa o mala fe del comerciante, sino a situaciones fuera de su control (como la pandemia), y no resulta justo castigarlo por ello. \\

\textbf{¿Qué efectos produce la quiebra de una empresa sobre el mercado?} &
El Estado deja de cobrar impuestos (Ingresos Brutos, Ganancias, IVA) y tiene más gastos (subsidio por desempleo); desaparece un bien o servicio del mercado y, con menos oferta y la misma demanda, suben los precios. Por eso aparece el salvataje. \\

\textbf{¿Qué es el salvataje (cramdown)?} &
Mecanismo por el cual se desinteresa al deudor, calculando cuánto valdría la empresa, y esta puede ser adquirida por un tercero con mayor capacidad y flujos, que la saque a flote y evite el desempleo y la caída del consumo. \\

\textbf{¿Qué diferencia hay entre un procedimiento privatista y uno publicista?} &
En el privatista, el deudor negocia directamente con los acreedores y el control lo ejerce una junta de acreedores; el juez, al homologar, solo verifica las formas. En el publicista, el control lo ejercen el juez y el síndico, y el juez no aprueba un plan que el síndico informa como inviable. \\

\midrule
\multicolumn{2}{p{0.93\textwidth}}{%
\textbf{Resumen:} Cada Estado puede concebir el procedimiento de insolvencia como castigo, como oportunidad de prevención o como instrumento de protección del mercado, y puede darle un tinte privatista o publicista según quién controle los acuerdos. Estas diferencias entre regímenes internos repercuten luego en el plano internacional.} \\
\end{longtable}

% ----------------------------------------------------------
\chapter[El caso Swift]{El caso Swift: multinacionales y descorrimiento del velo societario}
% ----------------------------------------------------------

\section{La multinacional de composición celular}

El caso Swift ilustra cómo un Estado puede controlar un acuerdo que formalmente cumple los requisitos legales. Swift era una empresa de estructura multinacional, de las que se dice que tienen una \textbf{composición celular}.

\begin{definition}{Multinacional de composición celular}
Estructura en la que existen vínculos entre empresas y que se diferencia de las relaciones entre controlantes y controladas, o de las que tienen participación, porque hay \textbf{un único plan, con una única cabeza que toma las decisiones}, y para ello esa cabeza se vale de distintas otras empresas.
\end{definition}

Se la llama «celular» porque todo se conecta con todo. Las empresas pueden vincularse de manera \textbf{vertical} o \textbf{horizontal}:

\begin{itemize}
    \item \textbf{Vertical:} se utilizan empresas de la misma área, todas vinculadas con la misma cadena de producción para llegar al mercado. Por ejemplo: una empresa ganadera, una de transporte ganadero, un matarife y un frigorífico.
    \item \textbf{Horizontal:} las empresas pueden tener diversidad, de modo que no se ve todo lo que hace una misma cabeza, y la conexión no resulta evidente aunque aparezcan como empresas distintas.
    % TODO: contenido incompleto o ambiguo en las notas originales (la explicación de la modalidad horizontal no está desarrollada)
\end{itemize}

En el caso Swift se aplicó el \textbf{descorrimiento del velo societario} para establecer que todas las sociedades estaban manejadas por una sola matriz: Deltec Argentina, que dependía de Deltec Internacional.

Desde la perspectiva de la persona jurídica, una multinacional es una única empresa y, por lo tanto, una única persona, con un solo patrimonio. Aunque las empresas hayan sido inscriptas como independientes, con un responsable asignado y un patrimonio propio, si se puede vincular que todas tienen la misma matriz y el mismo poder decisorio que va manejando cómo opera cada una, se está ante una multinacional de composición celular y, por lo tanto, \textbf{hay un único patrimonio}. Cada empresa sería como el riñón, los brazos o las piernas de una única persona jurídica: el mecanismo que utiliza un único cerebro, que toma las decisiones, para operar.

\begin{important}
Una multinacional de composición celular es \textbf{una sola persona jurídica con un único patrimonio}, aunque las sociedades estén inscriptas como independientes: el cerebro es uno solo y las empresas son sus órganos.
\end{important}

\section{El caso Swift: cómo se perjudicaba al accionista local}

El grupo Swift incluía, entre otras, a los frigoríficos La Negra y La Blanca y a una empresa de préstamos (una financiera). Swift se endeudaba y las empresas del grupo se conocían entre sí. Cuando se decretó la quiebra del frigorífico La Negra, el juez recibió el acuerdo con todas las conformidades aprobadas en asamblea, es decir, formalmente correcto. Un grupo de ganaderos formuló una objeción, denunciando una trama compleja.

La operatoria no se limitaba a un período previo: abarcaba \textbf{toda la actividad} del grupo.

\begin{itemize}
    \item Swift vendía a las empresas del grupo a precios más baratos que los de otros productos del mercado.
    \item Swift compraba a las empresas del grupo a precios superiores a los de plaza.
    \item Es decir, \textbf{compraba caro y vendía barato}: lo que perdía Swift lo ganaba el grupo empresario a través de las otras empresas.
    \item El perjudicado era \textbf{el accionista auténtico}, es decir, el inversor local, que había comprado acciones de esa empresa. Cuando la empresa quiebra, pierde él, porque lo que se perdía en un lugar se ganaba en otro: el grupo no pierde, pierde el inversor local.
\end{itemize}

El juez analizó todas estas cuestiones y \textbf{no homologó el acuerdo}, aplicando la teoría del descorrimiento del velo societario. La ley establecía que, si se obtenía la mayoría de acreedores y la mayoría de capital, el juez solo debía sumar; sin embargo, el juez puede mirar si el acuerdo es realmente viable o si existía algún problema.

Cuando se decretó la quiebra de Swift y de todo el grupo dentro de la República Argentina, fue inhabilitado por diez años, plazo que la ley establecía entonces como inhabilitación. Después se rehabilitó y volvió a ingresar al mercado.

\begin{important}
\textbf{Lección del caso Swift.} Aunque el acuerdo tenía las conformidades formales, el juez no homologó porque aplicó el descorrimiento del velo societario: el grupo ganaba lo que Swift perdía y el perjudicado era el accionista local. Muestra que el Estado, además de proteger el mercado y a los acreedores, puede introducir \textbf{mecanismos de contralor a cargo del propio Estado o a cargo de las partes interesadas}, dando al procedimiento un tinte más publicista o más privatista. En un procedimiento de control publicista, el juez puede examinar si el acuerdo es viable y legítimo, y no solo sumar mayorías.
\end{important}

\subsection*{Cuadro Cornell del capítulo}
\addcontentsline{toc}{section}{Cuadro Cornell del capítulo}

\begin{longtable}{>{\raggedright\arraybackslash}p{0.27\textwidth} >{\raggedright\arraybackslash}p{0.65\textwidth}}
\toprule
\textbf{Preguntas / palabras clave} & \textbf{Notas / respuestas} \\
\midrule
\endfirsthead
\toprule
\textbf{Preguntas / palabras clave} & \textbf{Notas / respuestas} \\
\midrule
\endhead
\bottomrule
\endlastfoot

\textbf{¿Qué es una multinacional de composición celular?} &
Un conjunto de empresas vinculadas, con un único plan y una única cabeza que decide, que se vale de distintas empresas (vinculadas de manera vertical u horizontal) para operar. Es una sola persona jurídica con un único patrimonio, aunque las sociedades figuren como independientes. \\

\textbf{¿Qué se entiende por descorrimiento del velo societario en el caso Swift?} &
Se aplicó para establecer que todas las sociedades del grupo estaban manejadas por una sola matriz (Deltec Argentina, dependiente de Deltec Internacional), de modo que había una única empresa y un único patrimonio. \\

\textbf{¿Cómo se perjudicaba al accionista local?} &
El grupo compraba caro y vendía barato dentro de sus propias empresas: lo que perdía Swift lo ganaba el grupo. Cuando la empresa quebraba, el grupo no perdía y perdía el inversor local, titular de acciones. \\

\textbf{¿Qué decidió el juez y por qué?} &
No homologó el acuerdo, aunque tenía las conformidades formales aprobadas en asamblea, porque aplicó el descorrimiento del velo societario: el juez puede examinar si el acuerdo es viable y no limitarse a sumar mayorías. \\

\textbf{¿Qué enseña el caso sobre el rol del Estado?} &
Que el Estado puede introducir mecanismos de contralor a cargo de sí mismo o de las partes interesadas, y así dar al procedimiento un tinte más publicista o más privatista. \\

\midrule
\multicolumn{2}{p{0.93\textwidth}}{%
\textbf{Resumen:} Un grupo multinacional de composición celular, con una única cabeza y un único patrimonio, perjudicó al accionista local mediante operaciones de compra cara y venta barata entre sus empresas. El juez, aplicando el descorrimiento del velo societario, se negó a homologar un acuerdo formalmente correcto.} \\
\end{longtable}

% ----------------------------------------------------------
\chapter[Criterios para la insolvencia internacional]{Criterios para regular la insolvencia internacional}
% ----------------------------------------------------------

\section{La insolvencia en un área de integración}

En un área de integración son muchas las variables respecto de la ley nacional que podría conducir la aplicación del procedimiento: por ejemplo, el domicilio del deudor o el establecimiento del deudor.

En un espacio en el que los Estados se comprometieron a la \textbf{igualdad de trato}, al \textbf{libre tránsito} y a permitir a las empresas extranjeras establecerse, la insolvencia plantea una situación distinta. Por ejemplo, un comerciante brasileño que contrató en la Argentina gracias al libre tránsito podría encontrarse con que, por ser acreedor extranjero, cobra después de que cobraron todos los acreedores nacionales.

Estas cuestiones hacen necesario, dentro de un área de integración económica, regular el procedimiento de insolvencia. Pero la cantidad de posibilidades hace muy difícil ponerse de acuerdo: es un tema controversial, según cómo regula cada Estado.

\section{Sistema extraterritorial (universal)}

Además de lo expuesto, el Estado puede adoptar un criterio de aplicación de la ley. La insolvencia puede ser internacional si se admite que el patrimonio está internacionalmente disperso, o puede ser \textbf{territorial}: en este último caso los efectos se dan solamente respecto de los bienes, créditos y deudas que tengan relación con el territorio de la República, aunque el patrimonio esté internacionalmente disperso.

\begin{definition}{Criterio extraterritorial}
Si el deudor tiene su domicilio en el país, la ley de ese Estado se aplica a todos los deudores, a todos los acreedores y a todos los bienes, sin importar dónde se encuentren. Basta determinar dónde tiene el deudor su domicilio comercial, el centro principal de los negocios o la sede (según las denominaciones), y ahí se aplica la ley, a todos. Los acreedores, sin importar el lugar de pago ni la nacionalidad, y los bienes, estén en el extranjero o en el país, forman \textbf{una sola masa}, tratada por un solo juez.
\end{definition}

Es el equivalente, en el ámbito internacional, del \textbf{fuero de atracción} que opera a nivel interno: así como no importa que el bien o el acreedor estén en Tucumán, en la provincia de Buenos Aires o en Neuquén, y un solo juez los trata, en el plano internacional la ley del juez competente se aplica extraterritorialmente.

\subsection{Fundamento 1: la persona es única y su patrimonio también}

Quiebra un comerciante o una empresa, que es una única persona. Lo que sabe el Derecho Internacional Privado es que corresponde aplicar la \textbf{ley personal}; si la ley personal es la del domicilio, se aplica la ley del domicilio.

\begin{itemize}
    \item Lo que se liquida no son créditos a favor del deudor y bienes aislados, sino \textbf{la totalidad del patrimonio}: se venden los bienes, se cobran los créditos y con el producido se cancelan las deudas.
    \item El patrimonio de la persona es único: es un atributo de la personalidad. (Se mencionó que los alemanes tienen la teoría del patrimonio de afectación, mientras que aquí se sostiene que el patrimonio es único.)
    \item Si lo que se liquida es un patrimonio, es lógico que se aplique una sola ley: la del titular de ese patrimonio, es decir, la del lugar donde tiene la sede principal.
\end{itemize}

\subsection{Fundamento 2: beneficia a los acreedores}

Los acreedores con crédito preferencial (hipotecarios, con garantía real de hipoteca o prenda) son acreedores privilegiados, con prioridad de pago respecto del bien que garantiza su crédito. El resto de los acreedores, los \textbf{quirografarios}, no tienen un bien que garantice la prestación: su garantía es la \textbf{prenda común de los acreedores}, es decir, el patrimonio en su totalidad.

Por lo tanto, conviene que todo se resuelva en un único procedimiento: cada cual verifica su crédito, se paga a los privilegiados y de lo que queda se reparte. Es el patrimonio en su totalidad lo que interesa, porque era la garantía de los acreedores.

Estos son los principales fundamentos de los países que adoptan el criterio de la extraterritorialidad para regular una quiebra internacional.

\section{Sistema territorialista}

Los países que adoptan un criterio territorialista aplican la ley a la totalidad de los bienes, créditos y deudas, pero solamente a los que están en el territorio del Estado.

\begin{definition}{Criterio territorialista}
Se debe probar la insolvencia (estado de cesación de pagos) conforme al criterio de la ley local, y la ley se aplica a todos los bienes, créditos y deudas situados dentro del territorio de la República. El estatus de insolvente \textbf{no surte efectos extraterritoriales}.
\end{definition}

El sistema territorialista no tiene como foco la persona, sino que el objetivo de la quiebra es pagar las deudas: se trata de reunir los bienes para ello, en un proceso económico y patrimonial importante. Si lo importante son los bienes, sean activo o pasivo, rige un principio de derecho internacional: a los bienes se les aplica la \textbf{ley del lugar de su situación}, así como a la persona se le aplica la ley personal. Por eso no se puede aplicar la ley argentina a un departamento situado en Uruguay.

\begin{itemize}
    \item \textbf{Juez competente:} el del lugar donde el deudor entró en estado de cesación de pagos, acreditado conforme a las características exigidas por la ley local.
    \item \textbf{Participantes:} todos los acreedores que tengan un lugar de cancelación de créditos en la República, y todos los deudores, no importa dónde vivan, que se hayan comprometido a pagar en la República. Por ejemplo, el extranjero que prometió pagar con un depósito bancario o con un cheque contra un banco con domicilio en el país: ese crédito se hace ejecutar en el territorio.
    \item \textbf{Alcance:} créditos, deudas y bienes que se encuentran en el territorio de la República; ahí se cierra. Aunque la quiebra sea internacional, no tiene efectos extraterritoriales.
\end{itemize}

\subsection{Fundamento: proteger a la empresa y evitar el «efecto dominó»}

\begin{itemize}
    \item El sistema no mira a la persona, sino a los bienes; su principal objetivo es cancelar las deudas, y se analiza solo una pequeña porción del patrimonio.
    \item Se favorece \textbf{la recuperación de las empresas internacionales}. A veces se entra en insolvencia no por ser un estafador ni un desastre, sino porque la situación de un mercado o de una economía en particular es desfavorable. Si la economía argentina resulta desfavorable a la empresa y esta entra en quiebra en la Argentina, extender esa declaración a Brasil, Uruguay y Paraguay haría desaparecer a la empresa de esos lugares, con retracción de la oferta en esos otros países.
    \item La dificultad de un mercado local se extendería como \textbf{fichas de dominó}: al caer la primera, se la extiende al ámbito internacional y se perjudica a acreedores y mercados que quizás están en buena situación de pago respecto de esa empresa.
\end{itemize}

\subsection{Consecuencia: pluralidad de procesos}

La única manera de cobrar respecto de los bienes de otro Estado es probar que en ese otro Estado también hay cesación de pagos o insolvencia, conforme a los requisitos de la ley local. El sistema territorialista lleva así a la \textbf{pluralidad de procesos}: hay que abrir tantos procesos de quiebra como mercados o países en los que se pueda probar que la empresa está insolvente. Si no se puede probar la insolvencia, no se puede abrir la quiebra, que se limita al patrimonio local. La declaración de insolvencia de un Estado surte efecto solamente territorial; si no se la puede probar del otro lado, se recurre a un juicio ordinario para cobrar. No es posible extender el estado de insolvencia.

\section{Dos miradas antagónicas: internacionalistas y comercialistas}

\begin{itemize}
    \item Los libros de \textbf{derecho internacional} sostienen normalmente que el mejor sistema es el \textbf{universal}, porque respeta la extraterritorialidad de la aplicación de la ley de quien declaró la insolvencia. Es el principio fundamental del derecho internacional.
    \item Los libros de \textbf{derecho comercial} reconocen el fuero de atracción en el ámbito interno, pero, como los profesores de derecho comercial asesoran en sus estudios jurídicos a la empresa, tienden a protegerla y a no extender la responsabilidad. Sostienen que lo mejor es el sistema \textbf{territorialista}, porque da a la empresa la oportunidad de restablecerse más rápido: si se hace quebrar a las otras, ya no tiene remedio; en cambio, si siguen operando, quizás las otras continúan operando bien y la empresa puede volver a abrirse en el país. Se facilita así la supervivencia de la empresa.
\end{itemize}

Se trata de dos criterios totalmente antagónicos, con fundamentos distintos.

\begin{longtable}{>{\raggedright\arraybackslash}p{0.19\textwidth} >{\raggedright\arraybackslash}p{0.36\textwidth} >{\raggedright\arraybackslash}p{0.35\textwidth}}
\caption{Sistema extraterritorial (universal) y sistema territorialista}\\
\toprule
\textbf{Aspecto} & \textbf{Extraterritorial (universal)} & \textbf{Territorialista} \\
\midrule
\endfirsthead
\toprule
\textbf{Aspecto} & \textbf{Extraterritorial (universal)} & \textbf{Territorialista} \\
\midrule
\endhead
\bottomrule
\endlastfoot
\textbf{Foco del sistema} & La persona (el deudor). & Los bienes y el pago de las deudas. \\
\textbf{Qué se liquida} & La totalidad del patrimonio, que es único (atributo de la personalidad). & Solo los bienes, créditos y deudas situados en el territorio del Estado. \\
\textbf{Ley aplicable} & Ley personal: la del domicilio o sede principal del deudor, para todo. & Ley del lugar de situación de los bienes (no se puede aplicar la ley argentina a un departamento en Uruguay). \\
\textbf{Juez competente} & Un único juez: el de la sede o domicilio (fuero de atracción internacional). & El juez del lugar donde se probó la cesación de pagos conforme a la ley local. \\
\textbf{Efectos de la declaración} & Surte efectos en todos lados. & No surte efectos extraterritoriales. \\
\textbf{Acreedores que participan} & Todos, sin importar nacionalidad ni lugar de pago. & Los que tengan lugar de cancelación del crédito en el territorio. \\
\textbf{Cantidad de procesos} & Un único procedimiento (con posibles variantes). & Pluralidad de procesos: uno por cada país donde se pruebe la insolvencia. \\
\textbf{Fundamento principal} & Unidad de la persona y del patrimonio; la garantía común de los acreedores es todo el patrimonio. & Proteger a la empresa: evitar el efecto dominó y facilitar su recuperación y supervivencia. \\
\textbf{Quién lo defiende} & Profesores de derecho internacional. & Profesores de derecho comercial, que asesoran a la empresa. \\
\end{longtable}

\section{Uniformidad o armonización: ¿qué conviene en un espacio de integración?}

Lo conveniente en el Mercosur, o en cualquier espacio de integración, sería contar con las mismas reglas que rigen dentro de un mismo Estado, es decir, establecer reglas \textbf{uniformes o armonizadas} respecto de este sistema. Lo ideal sería la \textbf{uniformidad}: por ejemplo, que el porcentaje o el tipo de procedimiento que un crédito laboral tiene en la Argentina sea exactamente igual al que tiene en Uruguay y en los demás Estados. Esto sería valioso porque al trabajador se le garantizó igualdad de trato y libre tránsito en materia laboral; si luego, al quebrar la empresa, recibiera un trato discriminado, esa garantía se vería frustrada. Lo ideal sería una única ley con parámetros uniformes, el mismo para todos los acreedores, todos los deudores y todos los bienes.

Sin embargo, dada la diversidad de regímenes, es muy difícil que todos los Estados se pongan de acuerdo en un único sistema con determinadas características.

\subsection{Las leyes de coyuntura}

Un obstáculo adicional es que, en todos los Estados, muchas veces se modifica una ley para evitar una cuestión económica transitoria.

\begin{example}{La crisis de 2001 y la prohibición de abrir quiebras}
Durante la crisis de 2001, cuando el país se encontraba en cesación de pagos, se prohibió, durante determinado tiempo, abrir juicios de quiebra. Es decir, \textbf{la mayor crisis económica que tuvo la República Argentina no produjo ni una sola quiebra}, porque la ley impedía abrir el juicio. La razón era que habilitar la apertura de juicios de quiebra provocaría una catarata de desempleos. La prohibición se prorrogó durante algún tiempo.
\end{example}

\begin{example}{La ley sucesoria}
En materia de sucesiones se modificó la ley sucesoria y se suprimió el impuesto para que una persona determinada pudiera hacer la sucesión sin pagar; después el impuesto se volvió a establecer.
% TODO: contenido incompleto o ambiguo en las notas originales (el pasaje contiene referencias a personas y términos inaudibles)
\end{example}

El Estado toma a veces medidas respecto de la ley de quiebras que pueden no ser muy lógicas y que sirven para tapar un agujero transitorio, como modificar un artículo para evitar que alguna empresa quiebre. Esto dificulta que se logre una ley única, con los mismos principios, los mismos resultados y las mismas características.

\begin{important}
Lo ideal en un espacio de integración sería una \textbf{ley única con parámetros uniformes}, pero la diversidad de regímenes y las leyes de coyuntura que los Estados modifican para tapar crisis transitorias hacen \textbf{muy difícil} el acuerdo; de ahí que se hable de reglas «uniformes o armonizadas».
\end{important}

\subsection*{Cuadro Cornell del capítulo}
\addcontentsline{toc}{section}{Cuadro Cornell del capítulo}

\begin{longtable}{>{\raggedright\arraybackslash}p{0.27\textwidth} >{\raggedright\arraybackslash}p{0.65\textwidth}}
\toprule
\textbf{Preguntas / palabras clave} & \textbf{Notas / respuestas} \\
\midrule
\endfirsthead
\toprule
\textbf{Preguntas / palabras clave} & \textbf{Notas / respuestas} \\
\midrule
\endhead
\bottomrule
\endlastfoot

\textbf{¿Qué problema plantea la insolvencia en un área de integración?} &
Habiéndose comprometido los Estados a la igualdad de trato y al libre tránsito, un acreedor extranjero podría cobrar después de todos los acreedores nacionales. De ahí la necesidad de regular el procedimiento, aunque las variables son muchas y el tema es controversial. \\

\textbf{¿Cómo funciona el sistema extraterritorial (universal)?} &
Si el deudor tiene domicilio o sede en el Estado, su ley se aplica a todos los acreedores, deudores y bienes, donde sea que se encuentren, y un solo juez los trata (fuero de atracción internacional). Se funda en que la persona y su patrimonio son únicos y en que el patrimonio total es la prenda común de los acreedores. \\

\textbf{¿Cómo funciona el sistema territorialista?} &
La ley se aplica solo a los bienes, créditos y deudas situados en el territorio, según la ley del lugar de situación; la declaración de insolvencia no surte efectos fuera de él. Lleva a la pluralidad de procesos y busca evitar el efecto dominó y favorecer la recuperación de la empresa. \\

\textbf{¿Qué consecuencia tiene el sistema territorialista para cobrar bienes situados en otro Estado?} &
Hay que probar la insolvencia en ese otro Estado conforme a su ley, y abrir tantos procesos como países en que se pueda probar. Si no se puede probar, se recurre a un juicio ordinario, porque no es posible extender el estado de insolvencia. \\

\textbf{¿Quién defiende cada sistema y por qué?} &
Los profesores de derecho internacional prefieren el universal, por respetar la extraterritorialidad. Los de derecho comercial, que asesoran a la empresa, prefieren el territorialista, que protege a la empresa y facilita su supervivencia. Son criterios antagónicos con fundamentos distintos. \\

\textbf{¿Por qué es difícil una ley uniforme para el Mercosur?} &
Lo ideal sería una única ley con parámetros uniformes (por ejemplo, igual trato al crédito laboral en cada Estado), pero hay mucha diversidad de regímenes y los Estados modifican leyes por razones de coyuntura (por ejemplo, la prohibición de abrir quiebras durante la crisis de 2001). \\

\midrule
\multicolumn{2}{p{0.93\textwidth}}{%
\textbf{Resumen:} Los Estados pueden regular la insolvencia internacional con un criterio extraterritorial o universal (ley y juez del domicilio para todo el patrimonio, único) o con un criterio territorialista (efectos limitados a los bienes situados en el territorio, con pluralidad de procesos). Defienden uno u otro los internacionalistas y los comercialistas, respectivamente. En un espacio de integración lo ideal sería la uniformidad, pero la diversidad y las leyes de coyuntura la dificultan.} \\
\end{longtable}

% ----------------------------------------------------------
\chapter[Régimen vigente en el Mercosur]{Régimen vigente en el Mercosur: Tratados de Montevideo y Código Bustamante}
% ----------------------------------------------------------

\section{Normas aplicables}

En el ámbito del Mercosur, la República Argentina se encuentra vinculada con los demás países por los siguientes instrumentos, todos con disposiciones sobre insolvencia:

\begin{longtable}{>{\raggedright\arraybackslash}p{0.38\textwidth} >{\raggedright\arraybackslash}p{0.52\textwidth}}
\caption{Instrumentos aplicables en el Mercosur}\\
\toprule
\textbf{Instrumento} & \textbf{Alcance indicado en los apuntes} \\
\midrule
\endfirsthead
\toprule
\textbf{Instrumento} & \textbf{Alcance indicado en los apuntes} \\
\midrule
\endhead
\bottomrule
\endlastfoot
Tratado de Derecho Comercial Internacional de Montevideo (1889) & Vincula a la Argentina con países asociados; contiene disposiciones sobre insolvencia. \\
Tratado de Derecho Comercial Terrestre Internacional de Montevideo (1940) & Vincula concretamente a la Argentina con Uruguay y con Paraguay; contiene disposiciones sobre insolvencia. \\
Código de Derecho Internacional Privado de Sánchez de Bustamante & Del que son parte Brasil y Venezuela; tiene muy pocos artículos sobre insolvencia (ocho o nueve). \\
\end{longtable}

Se trata de normas establecidas para un ámbito internacional y no para un ámbito regional. No pueden utilizarse en su totalidad, porque no cabe aplicar una parte de un tratado a un Estado y otra parte del tratado a otro. Los Tratados de Montevideo y el Código Bustamante adoptaron el criterio de la \textbf{extraterritorialidad de la insolvencia transfronteriza}.

\begin{note}
Los apuntes no citan números de artículo ni transcriben textos normativos; las referencias a los instrumentos se reproducen tal como fueron expresadas.
\end{note}

\subsection{Quiebra civil}

El Código Bustamante, a diferencia de los Tratados de Montevideo, que regulan solo la quiebra del derecho comercial, \textbf{también prevé la quiebra civil}, es decir, la insolvencia de quien no ejerce comercio pero se endeudó (por ejemplo, quien no pagó la cuenta del colegio de un hijo o la patente del automóvil). En los Tratados de Montevideo, la quiebra civil está regulada en el \textbf{Tratado de Derecho Procesal}, no en materia sustancial sino en el derecho procesal.
% TODO: contenido incompleto o ambiguo en las notas originales (se señala la diferencia sin indicar su fundamento)

\section{Extraterritorialidad: ¿juicio único o pluralidad de procedimientos?}

La construcción lógica es que la extraterritorialidad lleva a la tramitación de un único juicio, por el fuero de atracción. Sin embargo, no es la única manera de tramitar un juicio de insolvencia internacional: puede darse un \textbf{único juicio} o una \textbf{pluralidad de procedimientos}. Esto hace que muchos digan que el sistema es mixto.

Lo importante para hablar de extraterritorialidad es que la declaración de insolvencia \textbf{surta efectos fuera de las fronteras del Estado}: se declara con base en la persona, en el lugar de su sede o de su residencia, y esa insolvencia tiene efectos en todos lados. Puede tramitarse mediante un juicio de atracción o abriendo un procedimiento en cada Estado.

La diferencia con el sistema territorial es que \textbf{no hay que probar que ese otro Estado está en situación de cesación de pagos}, porque quebró el comerciante: no se está enfocado en el patrimonio, y alguien no puede ser insolvente en un lado y solvente en el otro.

\section{Juez competente, ley aplicable y actos aislados}

En el criterio extraterritorial basado en la persona, el juez competente es el del asiento de esa persona, de ese comerciante o de esa empresa: el \textbf{domicilio}. El Tratado de 1889 hablaba de domicilio; después surgieron otros criterios, como el lugar principal de producción económica o la administración. Probado el estado de insolvencia ante el juez del domicilio, este decreta la insolvencia y se aplica la ley del domicilio, sin importar dónde se cayó en la insolvencia.

Si se trata de una insolvencia internacional, el deudor debió haber realizado operaciones internacionales; si no, la situación se agota dentro de la frontera y es una quiebra nacional.

\subsection{Primera hipótesis: actos aislados}

Cuando el deudor realizó \textbf{actos aislados} en el ámbito internacional y no tiene sede en él (no tiene una sede permanente, por lo que su actividad internacional es esporádica), todos sus acreedores, todos sus deudores y todos sus bienes tramitan ante un único juez: el del domicilio. El juez competente es el del domicilio y la ley aplicable es la del domicilio del comerciante o de la empresa, a todo. Hay \textbf{un único juicio}.

La expresión «actos aislados» está en plural y se discutió mucho cuántos actos debe haber realizado una empresa para considerarse aislado y a partir de qué número se considera actividad permanente. Se trata de una \textbf{cuestión de análisis fáctico respecto de cada empresa}.

\begin{example}{Actos aislados según la dimensión de la empresa}
No es lo mismo un artesano que vende por internet desde su casa que una gran empresa multinacional: una empresa de gran tamaño puede haber realizado cien operaciones internacionales, pero si realiza un millón, son un porcentaje ínfimo; en cambio, quien vende diez mantitas tejidas en el año y envía seis al exterior tiene el mayor porcentaje de su actividad en el exterior. Cuántos actos hacen falta depende de cada situación en particular.
\end{example}

\section{Actividad permanente y sede independiente}

\subsection{Segunda hipótesis: actividad a través de una sede}

Distinto es el caso en que el deudor ha realizado actividades en otros Estados \textbf{a través de una sede}, es decir, mediante una actividad permanente. No basta con que tenga una sede: se requiere que \textbf{la sede haya actuado de manera totalmente independiente} de la casa matriz o casa del domicilio, tanto a nivel direccional como económico.

\begin{itemize}
    \item \textbf{Independencia direccional:} si la sucursal negocia un contrato y debe enviarlo a la casa matriz para que dé su aprobación, no es independiente: la decisión la genera el principal.
    \item \textbf{Independencia económica:} si para concretar un negocio la sucursal tiene que pedir los fondos a la casa principal, no es económicamente independiente.
\end{itemize}

Esto es una cuestión de prueba: no importa lo que digan los papeles. Quien contrata con la sede debe haber tratado con un sujeto que nunca dijo depender de la casa matriz para llevar adelante el negocio.

\begin{important}
Si hay una sede, y por lo tanto actividad permanente en un Estado distinto del domicilio del deudor, y esa sede es \textbf{independiente a nivel de dirección y a nivel económico}, existiendo acreedores locales, \textbf{los acreedores locales tienen la facultad de pedir la apertura del procedimiento en ese territorio}.
\end{important}

Es acreedor local el que tiene lugar de cancelación del crédito en ese Estado. Este procedimiento no lo realiza el juez del domicilio: interviene el juez del lugar de la sede y se aplica la ley de la sede.

El procedimiento sigue siendo extraterritorial, porque no hay que probar el estado de cesación de pagos de la sede: lo que se probó fue la insolvencia del deudor en el lugar del domicilio. Aunque la sede, con la que se hicieron negocios y que es totalmente independiente, esté pagando al término, igualmente puede pedirse la apertura. Si la sede fuera dependiente, es decir, sin autonomía, como un mero lugar de despacho, se entra en la hipótesis del juicio único.

\begin{longtable}{>{\raggedright\arraybackslash}p{0.20\textwidth} >{\raggedright\arraybackslash}p{0.34\textwidth} >{\raggedright\arraybackslash}p{0.36\textwidth}}
\caption{Actos aislados o sede dependiente y sede independiente (Tratados de Montevideo)}\\
\toprule
\textbf{Aspecto} & \textbf{Actos aislados o sede dependiente} & \textbf{Sede independiente} \\
\midrule
\endfirsthead
\toprule
\textbf{Aspecto} & \textbf{Actos aislados o sede dependiente} & \textbf{Sede independiente} \\
\midrule
\endhead
\bottomrule
\endlastfoot
\textbf{Tipo de actividad en el exterior} & Esporádica, sin sede permanente; o sede que depende de la casa matriz (mero lugar de despacho). & Actividad permanente a través de una sede, con poder de decisión y capital propio. \\
\textbf{Independencia requerida} & No hay: las decisiones y los fondos los provee la casa del domicilio. & Independencia a nivel direccional y económico; es una cuestión de prueba. \\
\textbf{Procedimiento} & Juicio único, bajo el fuero de atracción. & Posibilidad de procedimientos autónomos o locales, a pedido de acreedores locales. \\
\textbf{Juez competente} & El juez del domicilio del deudor. & El juez del lugar de la sede, para el procedimiento local. \\
\textbf{Ley aplicable} & La ley del domicilio, para todo. & La ley de la sede, para el procedimiento local. \\
\textbf{Quién puede pedir la apertura} & Se tramita en el domicilio. & Los acreedores con lugar de cancelación del crédito en ese territorio (acreedores locales). \\
\end{longtable}

\section{El procedimiento paso a paso: apertura, edictos y plazos}

\begin{enumerate}
    \item Probada la insolvencia (estado de cesación de pagos, tres balances, etc.), se acude al juez del domicilio del deudor y se pide la declaración de insolvencia.
    \item El juez abre el procedimiento de insolvencia (quiebra, concurso u otro) \textbf{conforme a la ley de su Estado}.
    \item El juez ordena la \textbf{publicación de edictos}, en todos los Estados donde el deudor haya realizado algún tipo de actividad (dependiente, independiente, habitual, no habitual) o tenga bienes. En 1889 era la única manera de enterarse de una quiebra en otro Estado.
    \item Finalizada la publicación, a partir del día siguiente a la última publicación, comienza a correr un plazo de \textbf{60 días} para que los acreedores locales del lugar donde el deudor tenga una sede independiente soliciten la apertura del procedimiento local.
\end{enumerate}

\begin{longtable}{>{\raggedright\arraybackslash}p{0.25\textwidth} >{\raggedright\arraybackslash}p{0.65\textwidth}}
\caption{Plazos en los Tratados de Montevideo}\\
\toprule
\textbf{Tratado} & \textbf{Plazos} \\
\midrule
\endfirsthead
\toprule
\textbf{Tratado} & \textbf{Plazos} \\
\midrule
\endhead
\bottomrule
\endlastfoot
\textbf{Tratado de 1889} & Publicación de edictos por 60 días. \\
\textbf{Tratado de 1940} & Publicación de edictos por 30 días. \\
\textbf{Ambos tratados} & Plazo de 60 días, desde el día siguiente a la última publicación, para que los acreedores locales de un lugar con sede independiente soliciten la apertura local. \\
\end{longtable}

En los lugares donde hay sede independiente, el plazo de 60 días corre desde la última publicación para que cualquier acreedor presente la solicitud de apertura del concurso o de la quiebra en ese territorio.

\section{Si vencen los 60 días y nadie pide la apertura local}

Si terminan los 60 días y ningún acreedor local pidió la apertura (o no había acreedores locales), \textbf{todo lo que está en ese país tramita ante el juicio del domicilio}: ese juicio, ese juez y esa ley. Esto muestra cuánto pueden cambiar las diferencias entre una ley y otra: si no se pide la apertura en el Estado de la sede y no se presenta nadie, todo debe ir al procedimiento del domicilio, ante un juez más lejano y, además, bajo otra ley.

\begin{important}
Los Tratados de Montevideo regularon, partiendo de la extraterritorialidad, la posibilidad de un \textbf{juicio único} o de \textbf{procedimientos plurales} en caso de que se abrieran procedimientos locales.
\begin{itemize}
    \item \textbf{Juicio único:} si se realizaron actos aislados, o si las sedes son dependientes del domicilio, todo queda bajo fuero de atracción ante el juez del domicilio.
    \item \textbf{Juicios autónomos o locales:} en el lugar donde se pruebe que la sede actuó de manera totalmente independiente a nivel directivo y económico, y haya acreedores locales.
    \item Si no hay acreedores locales no se puede pedir la apertura de la quiebra local, y se deberán liquidar todos los bienes y remitirlos al juicio del domicilio.
\end{itemize}
\end{important}

\section{Acreedores locales, prioridad de pago y remanente}

En el procedimiento local no existe un patrimonio totalmente independiente. Lo que se da es \textbf{prioridad de pago}: los acreedores locales tienen prioridad respecto de los bienes situados en el territorio. Si hay remanente, lo que sobra va al juicio del domicilio.

Esto demuestra que el sistema \textbf{no es totalmente territorial}: si lo fuera, no habría obligación de enviar los bienes sobrantes, ni podría declararse una quiebra si la sede era solvente. Razonamiento propio de la extraterritorialidad: nunca podría haber una quiebra solvente en este sistema, porque se demostró la insolvencia y es una única persona; si es fallido en un lado, tiene que serlo en todos.

La sede independiente es activo y pasivo: para tener acreedores, la sede debe pagar a personas en ese lugar. Al ser independiente económicamente, no tiene que pedir dinero a la casa matriz para pagar a sus acreedores, lo cual debe fijarlo la ley. Si, según la ley local, se pagó a los preferenciales y privilegiados y sobró dinero, ese dinero debe ir al juicio del domicilio, porque se trata siempre de \textbf{un patrimonio único}: los acreedores que se pagaron tuvieron preferencia para cobrar, pero no se trata de que el excedente se quede en el lugar.

\begin{note}
El Tratado de 1940 prevé el supuesto en que no se abre el procedimiento local y todo debe tramitarse en el juicio del domicilio: en ese caso, el juez competente no forma un único patrimonio con esos bienes, sino que debe abrir «\textbf{bolsitas}», patrimonios sobre los cuales tiene prioridad de pago. Los bienes que vinieron de la Argentina conceden preferencia de pago a los acreedores argentinos sobre los bienes situados en su territorio; esto ocurre si la sede era independiente, y no si era dependiente.
% TODO: contenido incompleto o ambiguo en las notas originales (la explicación del Tratado de 1940 aparece interrumpida y corregida en medio de la exposición)
\end{note}

\section{El síndico y la comunicación entre procedimientos: «no es mixto, es extraterritorial»}

El síndico nombrado en el juicio del domicilio tiene facultades para representar a la masa en todas las quiebras autónomas. Puede presentarse en representación de la quiebra principal, controlar que no se incluyan créditos indebidos y promover acciones contra la verificación de algunos créditos. Cuida, además, que lo que sobre vaya al juicio principal, que es la quiebra originaria.

Esto muestra que existe \textbf{comunicación entre los distintos procedimientos}: no hay un territorialismo cerrado. Por eso algunos dicen que el sistema es mixto. La posición expuesta es otra: \textbf{no es mixto, es extraterritorial}, porque todos los principios que adopta vuelven a la extraterritorialidad. Lo único que admite, procesalmente, es la posibilidad de abrir un único procedimiento o varios en determinadas situaciones: si la sede es independiente y existen acreedores locales.

Se destacó también que los Tratados de Montevideo, pese a ser antiguos, alcanzaron un alto nivel de perfección técnica.

\begin{important}
Los Tratados de Montevideo parten de la \textbf{extraterritorialidad}: la declaración de insolvencia, fundada en la persona, produce efectos en todos los Estados. Admiten un \textbf{juicio único} (actos aislados o sedes dependientes) o \textbf{juicios autónomos} (sede independiente y acreedores locales), con prioridad de pago a los acreedores locales, remanente al juicio del domicilio y facultades del síndico para representar a la quiebra principal. Si nadie pide la apertura local, todo se tramita ante el juez y bajo la ley del domicilio.
\end{important}

\subsection*{Cuadro Cornell del capítulo}
\addcontentsline{toc}{section}{Cuadro Cornell del capítulo}

\begin{longtable}{>{\raggedright\arraybackslash}p{0.27\textwidth} >{\raggedright\arraybackslash}p{0.65\textwidth}}
\toprule
\textbf{Preguntas / palabras clave} & \textbf{Notas / respuestas} \\
\midrule
\endfirsthead
\toprule
\textbf{Preguntas / palabras clave} & \textbf{Notas / respuestas} \\
\midrule
\endhead
\bottomrule
\endlastfoot

\textbf{¿Qué normas rigen la insolvencia entre la Argentina y los demás Estados del Mercosur?} &
El Tratado de Derecho Comercial Internacional de Montevideo (1889); el Tratado de Derecho Comercial Terrestre (1940), con Uruguay y Paraguay; y el Código Bustamante, del que son parte Brasil y Venezuela, que incluye también la quiebra civil. Todos adoptan la extraterritorialidad y fueron pensados para un ámbito internacional, no regional. \\

\textbf{¿Extraterritorialidad significa siempre un único juicio?} &
No. Lo esencial es que la declaración de insolvencia surta efectos fuera de las fronteras; puede tramitarse un único juicio o una pluralidad de procedimientos. A diferencia del sistema territorial, no hay que probar la cesación de pagos en el otro Estado. \\

\textbf{¿Qué juez y qué ley se aplican si el deudor realizó solo actos aislados?} &
El juez y la ley del domicilio del deudor, en un juicio único donde se concentran todos los acreedores, deudores y bienes. Cuántos actos son «aislados» es una cuestión fáctica (el artesano que vende por internet no es una gran multinacional). \\

\textbf{¿Cuándo corresponde un procedimiento local?} &
Cuando el deudor tiene una sede con actividad permanente, independiente a nivel directivo y económico, y hay acreedores locales (con lugar de cancelación del crédito en ese territorio). Interviene el juez del lugar de la sede y se aplica su ley. \\

\textbf{¿Cuáles son los plazos de edictos y de solicitud de apertura local?} &
Edictos por 60 días en el Tratado de 1889 y por 30 días en el de 1940. En ambos, desde el día siguiente a la última publicación corre un plazo de 60 días para que los acreedores locales pidan la apertura local. \\

\textbf{¿Qué ocurre si nadie pide la apertura local?} &
Todo se tramita en el juicio del domicilio, ante su juez y bajo su ley. \\

\textbf{¿Qué sucede con los bienes y acreedores si se abre un procedimiento local?} &
Los acreedores locales tienen prioridad de pago sobre los bienes situados en su territorio; el remanente va al juicio del domicilio, y el síndico de la quiebra principal puede representarla en las quiebras autónomas. El patrimonio sigue siendo único. \\

\textbf{¿Es mixto el sistema de los Tratados de Montevideo?} &
Según la posición expuesta, no: es extraterritorial, porque todos los principios vuelven a la extraterritorialidad, y solo se admite, procesalmente, tramitar un único procedimiento o varios en determinadas situaciones. \\

\midrule
\multicolumn{2}{p{0.93\textwidth}}{%
\textbf{Resumen:} En el Mercosur rigen los Tratados de Montevideo de 1889 y 1940 y el Código Bustamante, todos extraterritoriales. La regla es el juicio único ante el juez y con la ley del domicilio, salvo que exista una sede independiente a nivel directivo y económico y acreedores locales, caso en que se admiten procedimientos autónomos con prioridad de pago local, remanente al juicio principal y representación del síndico. Por eso el sistema no es mixto sino extraterritorial.} \\
\end{longtable}

% ----------------------------------------------------------
\chapter[Derecho comparado: Ley Modelo y Unión Europea]{Derecho comparado: Ley Modelo y Reglamento de la Unión Europea}
% ----------------------------------------------------------

\section{Procedimiento principal y procedimientos secundarios}

Los Tratados de Montevideo anticiparon la técnica de la \textbf{Ley Modelo de Insolvencia Transfronteriza}, que habla de un procedimiento principal y de procedimientos secundarios. Lo mismo hace la \textbf{Unión Europea}, a través del reglamento sobre insolvencia transfronteriza.

Un reglamento es una norma que tiene \textbf{aplicabilidad inmediata, supremacía y efecto directo} en el territorio de todos los Estados. El de Bruselas es del año 2000. En lugar de establecer que el juez competente del procedimiento principal es el del domicilio, dispone que el juez competente para la apertura y la declaración de la insolvencia es el del lugar del \textbf{centro principal de intereses} del deudor.

\section{El centro de intereses principales del deudor}

El reglamento no se refiere a dónde está la asamblea, dónde se reúne el consejo directivo, dónde está el domicilio efectivo, dónde está la planta de la fábrica ni cuál es el domicilio administrativo inscripto: utiliza el concepto de \textbf{centro de intereses principales del deudor}. El juez al que se le pide analiza si el vínculo que tiene con ese deudor puede calificarse como tal. Se aplica tanto a personas jurídicas como a personas físicas.

\begin{definition}{Centro de intereses principales}
Criterio \textbf{fáctico} para determinar el juez competente: se analiza dónde desarrolla efectivamente el deudor sus intereses principales y dónde están, en lugar de atender a la constitución de la sociedad o al domicilio. Por ejemplo, en la quiebra de un fabricante de alimentos, probablemente se considere que los intereses principales están donde produce, donde se toma la decisión.
\end{definition}

\begin{example}{Un fallo inglés}
Hubo un fallo de Inglaterra, cuando todavía pertenecía a la Unión Europea, en el que se declaró la insolvencia, dentro de la Unión Europea, de una empresa norteamericana, considerando competente al juez inglés. Estados Unidos no es parte de la Unión Europea. El juez consideró que el desarrollo, el centro de producción y la distribución hacia toda la Unión Europea estaban en Londres; por lo tanto, dentro de la Unión Europea el centro de intereses principales estaba en ese país, y aplicó el reglamento.
\end{example}

En consecuencia, el juez competente es el del lugar donde tiene el centro de intereses principales el deudor, y los otros procedimientos autónomos o independientes, que el reglamento llama \textbf{procedimientos secundarios}, se abren donde tenga un establecimiento.

\section{Establecimiento y procedimiento secundario}

La existencia de un establecimiento no implica automáticamente que sea el centro de intereses principales del deudor: el reglamento los separa. No distingue entre dependiente o independiente, sino que define como establecimiento el lugar donde el deudor ejerce una actividad \textbf{no transitoria}, a través de \textbf{recursos humanos y bienes}.

\begin{definition}{Establecimiento (Reglamento europeo)}
Aquel lugar donde el deudor ejerce una \textbf{actividad no transitoria}, a través de \textbf{recursos humanos y bienes}. No se exige «permanencia», que requeriría algún acta constitutiva de que se estará en el lugar por los próximos cien años: por ejemplo, una locación por cinco años ya es una actividad no transitoria. Tener recursos humanos significa tener personal contratado a cargo propio, y no a cargo de la casa matriz; y debe haber bienes en ese territorio.
\end{definition}

En el lugar donde se ubica un establecimiento puede abrirse el procedimiento \textbf{secundario}. Los Tratados de Montevideo hablaban de procedimiento «autónomo» o «independiente». No puede haber un procedimiento secundario si no hubiera uno principal. En este procedimiento tampoco hay que probar el estado de cesación de pagos: se lo puede abrir si existe un establecimiento.

En el establecimiento es competente el juez del lugar y se aplica la ley del lugar. En el centro de intereses principales del deudor se aplica la ley de ese lugar.

\begin{example}{Centro de intereses principales y establecimiento}
Si el centro de intereses principales está en Francia, se aplica la ley de quiebra de Francia; si el deudor tiene un establecimiento en Italia, el juez es el italiano y se aplica la ley de quiebra de Italia. La solución es igual a la de los Tratados de Montevideo.
\end{example}

\section{Armonía legislativa, no uniformidad}

Aunque haya una única ley con texto directo y aplicabilidad inmediata, \textbf{la solución sustantiva no es la misma en cada Estado}. El reglamento utilizó \textbf{la armonía legislativa, no la uniformidad}: se aplica una ley en un Estado y otra en otro, de modo que el acreedor que verifica en un lugar lo hace en condiciones distintas de las del que verifica en otro. Ello se debe a que las cuestiones analizadas al inicio dependen de las necesidades de cada Estado, y es muy difícil ponerse de acuerdo.

Lo que se acordó es la \textbf{obligatoriedad de aceptar las facultades del síndico del procedimiento principal} para intervenir en los procedimientos secundarios. Todo lo demás es prácticamente igual o presenta los mismos principios, aunque no igual.

Una diferencia señalada es que el reglamento prevé la comunicación por \textbf{videoconferencia}, mientras que el Tratado de Montevideo de 1889 no la contempla y la comunicación se realiza por exhortos u oficios, lo que se relaciona con la tecnología de la época. Los principios y la técnica, sin embargo, se tomaron del ámbito americano del siglo XIX.
% TODO: contenido incompleto o ambiguo en las notas originales (la referencia al «siglo» y a la comunicación por «exhortos u oficios» está reconstruida)

\section{El procedimiento secundario: solo liquidación y excepción}

El procedimiento secundario \textbf{solamente puede ser de liquidación}. Si se lo abre porque existe uno principal, puede no haber insolvencia en ese lugar, pero los bienes deben liquidarse y lo que sobre va al procedimiento principal. No puede habilitarse un concordato o un salvataje, aunque la ley local lo prevea: solo el procedimiento principal puede transferir la titularidad a otra empresa mediante un cramdown, cambiando el dueño y la titularidad. En el secundario únicamente se liquida.

El reglamento crea también la posibilidad de que se abra un \textbf{procedimiento secundario sin apertura de procedimiento principal}. En ese caso, el lugar donde el deudor tiene el establecimiento pasaría a ser el principal y actuaría como tal cuando la ley del lugar del centro de intereses principales del deudor no habilite la apertura (por ser menos favorable).

\begin{example}{Hipótesis expuesta para el procedimiento secundario que actúa como principal}
Si la ley de la actividad principal solo permitiera pedir la apertura del concurso al propio deudor o a los acreedores locales, y el deudor no tuviera acreedores locales, la ley del centro de intereses principales no permitiría abrir el juicio, aunque existan muchos acreedores. Entonces se podría acudir a un establecimiento, pedir la quiebra, probar el estado de cesación de pagos, y esa quiebra pasaría a actuar como la quiebra principal.
\end{example}

La explicación añadida es que, al aplicarse distintas leyes, no son veintisiete Estados sino un número reducido, y ya existen leyes más o menos similares.
% TODO: contenido incompleto o ambiguo en las notas originales (la parte final del pasaje se encuentra truncada)

\begin{longtable}{>{\raggedright\arraybackslash}p{0.19\textwidth} >{\raggedright\arraybackslash}p{0.34\textwidth} >{\raggedright\arraybackslash}p{0.37\textwidth}}
\caption{Tratados de Montevideo y Reglamento de la Unión Europea}\\
\toprule
\textbf{Aspecto} & \textbf{Tratados de Montevideo (1889 / 1940)} & \textbf{Reglamento de la Unión Europea (Bruselas, 2000)} \\
\midrule
\endfirsthead
\toprule
\textbf{Aspecto} & \textbf{Tratados de Montevideo (1889 / 1940)} & \textbf{Reglamento de la Unión Europea (Bruselas, 2000)} \\
\midrule
\endhead
\bottomrule
\endlastfoot
\textbf{Naturaleza} & Tratados internacionales antiguos, pensados para un ámbito internacional, no regional. & Reglamento: norma de aplicabilidad inmediata, supremacía y efecto directo en todos los Estados. \\
\textbf{Punto de partida} & Extraterritorialidad de la insolvencia. & Procedimiento principal y procedimientos secundarios. \\
\textbf{Juez del procedimiento principal} & Juez del domicilio del deudor. & Juez del lugar del centro de intereses principales del deudor (criterio fáctico). \\
\textbf{Otros procedimientos} & Procedimientos «autónomos» o «independientes» donde exista una sede independiente y acreedores locales. & Procedimientos «secundarios» donde exista un establecimiento (actividad no transitoria con recursos humanos y bienes). \\
\textbf{Ley aplicable} & En el procedimiento local, la ley del lugar de la sede. & Siempre la ley del lugar: la del centro de intereses principales en el principal; la del establecimiento en el secundario. \\
\textbf{Prueba de la cesación de pagos en el procedimiento local} & No hay que probarla en el Estado local. & No es necesario para el secundario, que se abre por la existencia de un establecimiento. \\
\textbf{Alcance del procedimiento local} & Prioridad de pago a acreedores locales; el remanente vuelve al juicio del domicilio. & Solo liquidación; lo que sobra va al principal. Sin concordato ni salvataje. \\
\textbf{Síndico} & El síndico de la quiebra principal puede representarla en las quiebras autónomas. & Obligatoriedad de aceptar las facultades del síndico del procedimiento principal. \\
\textbf{Comunicación entre procedimientos} & Por los medios de la época: no prevé la videoconferencia. & Por videoconferencia. \\
\textbf{Solución legislativa} & Cada Estado aplica su ley. & Armonía legislativa, no uniformidad. \\
\end{longtable}

\subsection*{Cuadro Cornell del capítulo}
\addcontentsline{toc}{section}{Cuadro Cornell del capítulo}

\begin{longtable}{>{\raggedright\arraybackslash}p{0.27\textwidth} >{\raggedright\arraybackslash}p{0.65\textwidth}}
\toprule
\textbf{Preguntas / palabras clave} & \textbf{Notas / respuestas} \\
\midrule
\endfirsthead
\toprule
\textbf{Preguntas / palabras clave} & \textbf{Notas / respuestas} \\
\midrule
\endhead
\bottomrule
\endlastfoot

\textbf{¿Qué es el reglamento europeo de insolvencia?} &
Una norma de aplicabilidad inmediata, supremacía y efecto directo en el territorio de todos los Estados. Como la Ley Modelo, habla de procedimiento principal y procedimientos secundarios. \\

\textbf{¿Qué es el centro de intereses principales?} &
El lugar donde el deudor desarrolla efectivamente sus intereses principales (criterio fáctico), más allá del domicilio inscripto o del lugar donde se reúne la asamblea. Determina el juez competente para el procedimiento principal, y se aplica su ley. \\

\textbf{¿Qué es un establecimiento?} &
El lugar donde el deudor ejerce una actividad no transitoria, a través de recursos humanos y bienes. No se exige permanencia. Permite abrir un procedimiento secundario, con juez y ley del lugar. \\

\textbf{¿En qué se parecen los procedimientos secundarios y los autónomos de los Tratados de Montevideo?} &
En ambos casos no hay que probar la cesación de pagos en el lugar y se aplican el juez y la ley del lugar de la sede o del establecimiento. La solución es igual a la de los Tratados de Montevideo. \\

\textbf{¿Armonía legislativa o uniformidad?} &
El reglamento utiliza la armonía legislativa, no la uniformidad: cada Estado aplica su ley y la solución sustantiva no es la misma. Lo obligatorio es aceptar las facultades del síndico del procedimiento principal en los procedimientos secundarios. \\

\textbf{¿Qué limita al procedimiento secundario?} &
Solo puede ser de liquidación (sin concordato ni salvataje) y lo que sobra va al principal. Excepcionalmente puede abrirse sin procedimiento principal y actuar como principal cuando la ley del centro de intereses principales no habilite la apertura. \\

\midrule
\multicolumn{2}{p{0.93\textwidth}}{%
\textbf{Resumen:} El reglamento europeo y la Ley Modelo ofrecen una solución equivalente a la de los Tratados de Montevideo con otra técnica: procedimiento principal ante el juez del centro de intereses principales del deudor y procedimientos secundarios, solo de liquidación, en cada establecimiento, con armonía legislativa y no uniformidad.} \\
\end{longtable}

\end{document}
